
\section{Why Does the Lutheran Free Church Attempt to Carry On Home Mission Work?}

“Folkebladet,” April 26, 1916

E. B.

To this question there can no doubt be various answers, since the question may be viewed and judged from several sides.

The first answer, and the one which seems to me the most immediate, is this: The Lutheran Free Church attempts to carry on home mission work—as well as foreign mission—because it must.

When that which we call Christianity, or perhaps more properly life in God, has become—not a heavy burden which we laboriously drag along with us—but our happiness, the hidden source of our life’s joy, then it is, both for the individual and for the congregation, a natural inward need to labor that others also may be brought to share in this life and this happiness.

It is also both right and natural that in this work we should think first and foremost of our own people. The first Christian missionaries were commanded to begin at home, “beginning from Jerusalem.”

Life in God Jesus himself has described as the greatest happiness, indeed as blessedness. And he has compared it to a treasure which is worth so much that everything ought to be sold in order that one may come into possession of it.

And this treasure is unlike earthly treasures. These—if we obtain any of them—we would rather keep for ourselves. But not this treasure. It is sufficient, it is great enough not merely for some of us, but for all, all.

And it has this peculiar quality, that the more we give out of it to others, the more we ourselves have remaining. Never is our peace with God deeper, our joy, our life in him richer than when, perhaps without sparing ourselves, we labor to help our brothers and sisters find the “treasure” and come to possess it.

Perhaps one or another of you who read these words has experienced that it is so.

Or was it not so, that on the days when you most earnestly and tirelessly used your strength in your Master’s service, resolved to go on his errand, then you knew in your heart something of that peace which surpasses all understanding, and which Jesus promises to his own?

And on the other hand, when you “spared yourself,” were unwilling to do anything for the Lord, and neglected your opportunities, did you not then experience that the heart became troubled and the soul without peace?

And as it is with the individual believer, so also with the congregation. The congregation which concerns itself only with itself and its own “edification” will languish and lose its strength and the blessing of God.

To the question why we attempt to carry on home mission work, I believe that most of those for whom life in God has become a blessed reality will answer: We carry on home mission because we must.

This work has become an impelling necessity for us, insofar as we are from time to time permitted to believe that we ourselves belong to the Lord.

Paul speaks of how “the love of Christ compels.”

God grant us much of this love which “compels.” Then also the work which we call our home mission will be done with joy and willingness.
